Use the Hauptmodul normalized by the values $0,\infty,-16/27$
at the CM points of discriminants $-4,-3,-24$.
In a common weighted-projective representative,
\cite[equations (2) and (14), pp.~3 and 17]{linyang} gives
\[
 J_2=12(j+1),\quad J_4=6(j^2+j+1),\quad J_{10}=j^3,
 \qquad \psi_4=J_2^2-24J_4=144j,\quad \chi_{10}=-j^3/4.
\]
The conventions of \cite[\S2.2, p.~4]{lny} are
$\psi_4=2^{-2}\sum_{m\ {\rm even}}z_m^8$ and
$\chi_{10}=-2^{-14}\prod_{m\ {\rm even}}z_m^2$.
The quotient $\chi_{10}^2/\psi_4^5$ is unchanged by weighted
rescaling. Since our $j$ is finite and nonzero, its denominator
is nonzero, and substitution gives
\[
 j=16\cdot144^5\frac{\chi_{10}^2}{\psi_4^5}
  =3779136\frac{(\prod_{m\ {\rm even}}z_m)^4}
                 {(\sum_{m\ {\rm even}}z_m^8)^5}.
\]
The two homogeneous polynomials have degree $40$ and coefficient
$\ell^1$ norms $3779136$ and $10^5$. Hence
$h(j)\le40h([z_m])+22$.
The product and sum use the ten even theta nulls of the same
principally polarized surface. Adjoining the six zero odd
coordinates does not change their projective height.

In \cite{pazuki}, the level $r=2$ theta coordinates are the sixteen
half-characteristic values used here; the six odd values vanish.
The permitted changes of theta structure permute coordinates and
multiply them by roots of unity. His Euclidean projective height
dominates our supremum projective height. Put
$X=\max(1,h_F(A))$ and
$X_D=\max(1,h_F(A)+\log\pi)\le3X$.
His Corollary~1.3(2), in the Deligne metric, gives
\[
 h_\Theta(A)\le\tfrac12 X_D+C_2\log(X_D+2),
 \qquad C_2=1000\cdot2^4(\log(2^4))^5\le1000\cdot2^{24}.
\]
Using $\log(X_D+2)\le2X_D$ proves
$h([z_m])\le10000\cdot2^{24}X$.

Write $j=a/b$ in lowest terms. The nonzero integer polynomial
$-57600bF(X)^2-aG(X)^6$ has $u$ as a root and coefficient norm at
most $2\cdot4000^6\max(|a|,|b|)$. Remove any initial power of $X$
and use the actual highest nonzero coefficient. The rational-root
theorem bounds both numerator and denominator of $u$, giving
$h(u)\le h(j)+\log(2\cdot4000^6)<h(j)+74$.

A field extension for the theta structure changes neither absolute
height. Since $40\cdot10000\cdot2^{24}+96<10^{13}$, the preceding
inequalities give
\[
 h(u)\le10^{13}\max(1,h_F(A))
 \le10^{13}\max(1,C_{\rm cm},Q_R^{17})M_S(u)^{24}
 =\tfrac12 H_6(S;R)M_S(u)^{24}.
\]
Every estimate used only the displayed lower bound on $R$, so the
conclusion holds for every permitted larger $R$ as well.
\end{proof}

